\section{Evaluaci\'on}
\label{sec:evaluation}

\subsection{Metodolog\'ia}

Los datos de referencia de cada una de las tres familias se
produjeron mediante ingenier\'ia inversa manual, documentados como un
informe de analista completo y nunca derivados del propio resultado
de la pipeline. Esta independencia es una salvaguarda metodol\'ogica
deliberada, no un detalle incidental: ambos motores dependen de
herramientas de terceros (YARA, FLOSS, firmas de la comunidad de CAPE)
que pueden por s\'i mismas mapear mal una t\'ecnica, de modo que los
datos de referencia tuvieron que establecerse sin consultar los propios
artefactos extra\'idos por la pipeline, ya que de lo contrario validar
la pipeline frente a ella ser\'ia circular. En concreto, el mismo
protocolo se aplic\'o a las tres familias, en este orden, y solo el
\'ultimo paso llega a tocar el resultado propio de la pipeline:

\begin{enumerate}
  \item \textbf{Desensamblado est\'atico independiente.} Cada muestra
    se carg\'o en Ghidra y se descompil\'o funci\'on por funci\'on, cada
    funci\'on de inter\'es renombrada para reflejar su prop\'osito
    establecido, sin consultar la propia extracci\'on de importaciones
    o cadenas de la pipeline.
  \item \textbf{Observaci\'on din\'amica independiente.} La detonaci\'on
    en sandbox se observ\'o directamente (secuencia de llamadas a API,
    ficheros soltados, actividad de red), registrando el comportamiento
    a partir de la traza en bruto en lugar de a trav\'es de las propias
    etiquetas de firma de la comunidad de CAPE, que \S\ref{sec:dynamic}
    documenta como necesitadas de curaci\'on por s\'i mismas.
  \item \textbf{Recuperaci\'on independiente de configuraci\'on y
    claves.} Cuando una muestra cifraba u ofuscaba su configuraci\'on,
    la clave o el algoritmo se recuper\'o manualmente, completamente
    separado de los propios m\'odulos automatizados de fuerza bruta y
    recuperaci\'on XOR de la pipeline.
  \item \textbf{Solo entonces, la comparaci\'on.} Se ejecut\'o la
    pipeline y se compar\'o su \code{attck\_mapping.json} con la tabla
    de T\'ecnicas ATT\&CK del informe terminado y producido de forma
    independiente (la auditor\'ia de falsos positivos,
    \S\ref{sec:fp-audit}). Por construcci\'on, cada discrepancia que
    encuentra es un desacuerdo entre dos artefactos producidos de forma
    independiente, no la pipeline comprobada frente a sus propias
    suposiciones.
\end{enumerate}

La Tabla~\ref{tab:validation} da un ejemplo concreto por familia de
evidencia que este protocolo recuper\'o y que la pipeline automatizada
nunca toca, que es lo que convierte cada informe en una comprobaci\'on
independiente y no en una reformulaci\'on de lo que la herramienta ya
asume.

\begin{table}[h]
\centering
\small
\begin{tabular}{@{}lp{9.5cm}@{}}
\toprule
\textbf{Muestra} & \textbf{Recuperado de forma independiente, la pipeline no lo toca} \\
\midrule
Akira & Tabla ATT\&CK de 17 filas construida a partir de descompilaci\'on de Ghidra a nivel de funci\'on de los mecanismos internos de threading, logging y cifrado de la muestra. \\
\rowline
RoningLoader & Los cuatro componentes de la cadena (instalador, \code{D3D11InstallHelper.dll}, driver rootkit, cliente RAT) descompilados por separado, incluyendo identificar \code{DoD3D11InstallUsingMSI} del auxiliar como una reimplementaci\'on compilada en Rust de \code{env::var()}. \\
\rowline
WhiteSnakeStealer & Clave RC4 recuperada mediante an\'alisis de texto plano conocido, y el formato completo del paquete C2 \code{WSR\$} (XML de informaci\'on del sistema cifrado con RC4 y comprimido con Deflate, con una clave envuelta en RSA) reconstruido directamente a partir del desensamblado. \\
\bottomrule
\end{tabular}
\caption{Un ejemplo concreto por familia de evidencia recuperada de
forma independiente para construir los datos de referencia, aplicando
el protocolo de cuatro pasos anterior.}
\label{tab:validation}
\end{table}

La tabla de T\'ecnicas ATT\&CK de cada informe se analiza despu\'es para
producir un fichero de datos de referencia estructurado
(\code{extract\_ground\_truth.py}), que el arn\'es de evaluaci\'on
compara con el propio \code{attck\_mapping.json} de la pipeline.

Se reportan dos modos de emparejamiento. El emparejamiento
\emph{estricto} exige una coincidencia exacta de identificador de
t\'ecnica. El emparejamiento a \emph{nivel de familia} trata una
t\'ecnica padre y cualquiera de sus subt\'ecnicas como equivalentes (un
hallazgo de la pipeline de \technique{T1547.001} cuenta frente a una
entrada de datos de referencia de \technique{T1547} y viceversa), bajo
la premisa de que un analista de seguridad que decide qu\'e hacer ante
un mecanismo de persistencia se preocupa primero por \emph{qu\'e}
familia de t\'actica y t\'ecnica est\'a presente, siendo la
granularidad de subt\'ecnica un refinamiento y no un hallazgo
distinto. Se calculan precisi\'on, exhaustividad y F1 bajo ambos modos
para cada muestra; la Tabla~\ref{tab:results} reporta las cifras a
nivel de familia, siendo las cifras m\'as estrictas materialmente m\'as
bajas solo donde realmente no se distingui\'o una subt\'ecnica
concreta.

\subsection{Resultados}

\begin{table}[h]
\centering
\small
\begin{tabular}{@{}lcccc@{}}
\toprule
\textbf{Muestra} & \textbf{Precisi\'on} & \textbf{Exhaustividad} & \textbf{F1} & \textbf{Notas} \\
\midrule
Akira (ransomware) & 0.92 & 1.00 & 0.96 & binario de una sola etapa \\
\rowline
WhiteSnakeStealer (.NET) & 1.00 & 0.80 & 0.89 & binario de una sola etapa \\
\rowline
RoningLoader (solo padre) & 0.85 & 0.68 & 0.76 & solo el cargador \\
\rowline
RoningLoader (+reenv\'io) & 0.84 & 0.84 & 0.84 & cargador + 3 componentes soltados \\
\bottomrule
\end{tabular}
\caption{Precisi\'on/exhaustividad/F1 a nivel de familia por muestra
validada.}
\label{tab:results}
\end{table}

Las dos filas de RoningLoader aislan exactamente lo que aporta el
bucle de reenv\'io (\S\ref{sec:resubmission}): puntuar \'unicamente el
binario cargador frente a unos datos de referencia escritos para
\emph{toda} su cadena de infecci\'on subestima la exhaustividad por
construcci\'on, ya que las t\'ecnicas propias del driver rootkit y del
cliente RAT son hallazgos reales que los datos de referencia esperan
pero que el binario padre por s\'i solo no puede presentar. La
exhaustividad sube 16 puntos en cuanto se analizan de forma
independiente los tres componentes soltados y se cuentan sus
hallazgos, a un coste de 1 punto de precisi\'on por la superficie
a\~nadida, una compensaci\'on que la mejora del F1 agregado
(0.76~$\rightarrow$~0.84) demuestra que mereci\'o la pena.

La precisi\'on es alta en las cuatro filas (0.84--1.00), es decir,
el peque\~no n\'umero de hallazgos que reporta la pipeline rara vez son
err\'oneos. Esta es la consecuencia directa e intencionada de la
filosof\'ia de combinaci\'on coherente de \S\ref{sec:static-rules} y la
disciplina de curaci\'on de firmas de \S\ref{sec:dynamic}: ambas
sacrifican algo de exhaustividad a cambio de alta confianza en que lo
que \emph{se} reporta es fiable, algo que la auditor\'ia de falsos
positivos a continuaci\'on existe para comprobar, no para asumir.

\subsection{La auditor\'ia de falsos positivos}
\label{sec:fp-audit}

La precisi\'on por s\'i sola no distingue ``la pipeline est\'a bien
calibrada'' de ``los datos de referencia resultan ser permisivos''. A
diferencia de un agregador de reputaci\'on que muestra la etiqueta de
cada motor o sandbox tal cual, cada discrepancia entre el resultado
de la pipeline y los datos de referencia encontrados en las tres muestras,
15 en total (5 Akira, 4 WhiteSnakeStealer, 6 RoningLoader), se
rastre\'o individualmente hasta una causa concreta y verificable en
lugar de aceptarse o descartarse solo por la cifra agregada. Cada una
se contrast\'o con el \emph{texto completo} del informe manual
correspondiente, no solo con su tabla resumen, ya que una t\'ecnica
puede estar gen\'uinamente discutida en la narrativa de un informe sin
aparecer en su propia tabla ATT\&CK.

El resultado de esta auditor\'ia, por causa:

\begin{itemize}
  \item \textbf{Diez errores gen\'uinos de la pipeline}, cada uno
    corregido en su causa ra\'iz: tres mapeos de regla est\'atica que no
    correspond\'ian con la definici\'on real de MITRE para la t\'ecnica
    afirmada (p.\ ej.\ \code{GetSystemInfo} mapeada a Descubrimiento de
    Informaci\'on del Sistema, \technique{T1082}, cuando la propia
    definici\'on de MITRE trata espec\'ificamente sobre detecci\'on de
    \emph{virtualizaci\'on/sandbox}, \technique{T1497}); un hueco en una
    comprobaci\'on de combinaci\'on que permit\'ia que una \'unica
    importaci\'on gen\'erica activara \technique{T1055} por s\'i sola; y
    seis mapeos de firmas de la comunidad de CAPE cuyo propio campo de
    evidencia en bruto, al inspeccionarlo, nombraba un objetivo
    inconsistente con la t\'ecnica afirmada (detallado en
    \S\ref{sec:dynamic}).
  \item \textbf{Dos correcciones de calibraci\'on de confianza}, en
    Akira y RoningLoader: el mismo hueco de comprobaci\'on de
    combinaci\'on de la correcci\'on de \technique{T1055} anterior
    reapareci\'o para \technique{T1070} (Indicator Removal) --
    \code{CreateFileW}/\code{CreateFile}, ya marcada como ``gen\'erica''
    en la tabla de mapeo, segu\'ia cont\'andose ciegamente como
    corroboraci\'on junto a \code{DeleteFileW}/\code{DeleteFile} y se
    promov\'ia a confianza \emph{alta}. Migrar los datos de referencia a
    ATT\&CK v19 (\S\ref{sec:limitations}) es lo que lo sac\'o a la luz:
    la antigua entrada de datos de referencia \technique{T1070.001}
    hab\'ia estado emparej\'andose con este hallazgo por pura
    coincidencia de prefijo de identificador bajo el emparejamiento a
    nivel de familia, no porque la evidencia tratara gen\'uinamente
    sobre el asunto real de esa entrada (el borrado de registros de
    eventos); mover esa entrada a su identidad v19
    (\technique{T1685.005}, que ya no comparte el prefijo
    \technique{T1070}) rompi\'o la coincidencia casual. Corregido de la
    misma manera que el caso de \technique{T1055}, excluyendo la
    importaci\'on/cadena gen\'erica de la corroboraci\'on en ambas
    direcciones. A diferencia de aquel caso, la evidencia subyacente
    aqu\'i es real (ambos binarios importan efectivamente estas
    funciones), simplemente no lo bastante espec\'ifica para confirmar
    la t\'ecnica afirmada, as\'i que el hallazgo permanece presente
    tras la correcci\'on, ahora con su confianza honesta
    \emph{media}/\emph{baja}, en lugar de desaparecer del todo -- una
    ilustraci\'on directa de la relaci\'on entre confianza y precisi\'on
    examinada en \S\ref{sec:confidence-calibration}, no un
    contraejemplo de ella.
  \item \textbf{Un hueco de extracci\'on de datos de referencia}, no un
    error de la pipeline: el propio informe manual de
    WhiteSnakeStealer discute la carga reflexiva de DLL en su texto
    narrativo, pero la fila correspondiente de \technique{T1620} faltaba
    en la tabla ATT\&CK del propio informe, de modo que la extracci\'on
    autom\'atica de datos de referencia nunca la captur\'o. El hallazgo
    de la pipeline era correcto; se corrigi\'o el documento fuente en
    lugar de la pipeline.
  \item \textbf{Dos dejados abiertos deliberadamente}, en RoningLoader
    (\technique{T1027}, \technique{T1497}): la evidencia de apoyo de la
    pipeline para ambos es real, pero no existe confirmaci\'on expl\'icita
    del analista en ning\'un lugar del texto del informe, a diferencia
    del caso de \technique{T1620} anterior. Editar los datos de
    referencia solo con la evidencia propia de la pipeline, sin
    confirmaci\'on textual independiente, ser\'ia circular. Ambos
    permanecen como hallazgos abiertos en lugar de resolverse en un
    sentido u otro.
\end{itemize}

Esta asimetr\'ia, un hueco de datos de referencia corregido y dos
discrepancias de aspecto similar dejadas abiertas, es deliberada: el
factor distintivo es si existe confirmaci\'on independiente en el
documento fuente, no si la evidencia propia de la pipeline parece
convincente. El resultado pr\'actico de ejecutar esta auditor\'ia se
aprecia en la columna de corregidos: 12 de las 15 discrepancias se
tradujeron en mejoras concretas y verificables al conjunto de reglas
(\S\ref{sec:coverage}) y a la l\'ogica de agregaci\'on
(\S\ref{sec:case-installer}) en lugar de absorberse en silencio dentro
de una cifra agregada, precisamente el paso diagn\'ostico que una
visualizaci\'on de veredictos uno junto a otro no puede ofrecer.

\subsection{Calibraci\'on de confianza}
\label{sec:confidence-calibration}

Si el modelo de reconciliaci\'on de confianza (\S\ref{sec:reconciliation})
es significativo y no meramente cosm\'etico, la precisi\'on deber\'ia
ser medible y mayor entre los hallazgos de confianza \emph{alta}
(aquellos corroborados entre fuentes o entre niveles de t\'ecnica) que
entre los hallazgos de \emph{media}/\emph{baja} de una sola fuente. En
las tres muestras validadas, 14 de las 15 discrepancias identificadas en
la auditor\'ia anterior tuvieron su origen en un hallazgo de una sola
fuente. La \'unica excepci\'on, \technique{T1497} de RoningLoader,
alcanz\'o confianza \emph{alta} mediante corroboraci\'on gen\'uina entre
fuentes (una combinaci\'on coherente de APIs anti-sandbox en la
evidencia est\'atica, tres firmas independientes de CAPE en la evidencia
din\'amica), y sin embargo sigue siendo uno de los dos hallazgos
dejados abiertos deliberadamente m\'as arriba, precisamente porque no
existe confirmaci\'on textual independiente en el informe manual, no
porque la evidencia de ninguna de las dos fuentes sea d\'ebil. Le\'ido
as\'i, esto no contradice el argumento de calibraci\'on: el modelo
reconoci\'o correctamente una evidencia fuerte y corroborada; lo que es
gen\'uinamente incierto es si los datos de referencia est\'an completos,
la misma asimetr\'ia que establece el caso de \technique{T1620} en otro
punto de esta auditor\'ia, no si el nivel de confianza asignado estaba
justificado. Esto es consistente con que el modelo de
reconciliaci\'on realiza una calibraci\'on real y no simplemente
reetiqueta hallazgos, aunque el tama\~no reducido de la muestra (tres
familias) implica que esto debe leerse como una se\~nal direccional, no
como una afirmaci\'on con potencia estad\'istica, una limitaci\'on
enunciada expl\'icitamente en \S\ref{sec:limitations}.
